%% kapa: exam=1 round=31 year=2020 period=1 form=A subject=economics subject_ko=경제학원론 seq=22 number=62
%% source: Q-Net 감정평가사 기출문제 (cst00309) artlSeq=5207219 file=extracted/1008/5207219/A형/A형/2020년도 제31회 감정평가사 1차 시험_1교시_A형.pdf
%% conversion: kordoc 4.24.51 -> tools/gen1.py (auto); review=proofread-sonnet; flags=none
\begin{question}{62}
효율임금이론에 관한 설명으로 옳지 않은 것은?
\begin{choices}
  \item 높은 임금을 지급할수록 노동자 생산성이 높아진다.
  \item 높은 임금은 이직률을 낮출 수 있다.
  \item 높은 임금은 노동자의 도덕적 해이 가능성을 낮출 수 있다.
  \item 효율임금은 시장균형임금보다 높다.
  \item 기업이 임금을 낮출 경우 생산성이 낮은 노동자보다 높은 노동자가 기업에 남을 확률이 높다.
\end{choices}
\end{question}
\answer{5}
